\documentclass{beamer}

\usetheme{Madrid}
\usecolortheme{seahorse}

\usepackage{graphicx}

\title{Sistemi Operativi}
\subtitle{Concetti base e architettura}
\author{Prof. Fedeli Massimo - Fabbrica Digitale}
\date{}

\begin{document}
	
	\frame{\titlepage}
	
	%------------------------------------------------
	\begin{frame}
		\frametitle{Sommario}
		
		\begin{itemize}
			\item Funzionamento di un calcolatore
			\item Obiettivi e funzioni di un sistema operativo
			\item Implementazione e avvio del sistema operativo
			\item Concetti di base sui sistemi operativi
			\item Servizi di un sistema operativo
			\item Funzionamento event-driven e dual mode
			\item Servizi e chiamate di sistema
			\item Programmi di sistema
			\item Struttura di un sistema operativo
		\end{itemize}
		
	\end{frame}
	
	%------------------------------------------------
	\section{Funzionamento di un calcolatore}
	
	\begin{frame}
		\frametitle{Funzionamento di un calcolatore}
		
		\begin{itemize}
			\item Un calcolatore è composto da una CPU e diversi controllori di dispositivi
			\item Comunicano tramite un bus
			\item Operano in modo concorrente contendendo l’accesso alla memoria centrale
		\end{itemize}
		
		\begin{center}
			\includegraphics[width=0.7\textwidth]{figures/calcolatore_architettura}
		\end{center}
		
	\end{frame}
	
	%------------------------------------------------
	
	\begin{frame}
		\frametitle{I/O e trasferimento dati}
		
		\begin{itemize}
			\item L'I/O avviene tra il dispositivo e il buffer locale del controller
			\item La CPU guida poi lo spostamento dei dati sul bus
			\item I dati vengono trasferiti tra buffer dei controller e memoria centrale
		\end{itemize}
		
		\begin{center}
			\includegraphics[width=0.7\textwidth]{figures/io_controller}
		\end{center}
		
	\end{frame}
	
	%------------------------------------------------
	
	\begin{frame}
		\frametitle{Organizzazione della memoria}
		
		\begin{center}
			\includegraphics[width=0.7\textwidth]{figures/memory_hierarchy}
		\end{center}
		
	\end{frame}
	
	%------------------------------------------------
	
	\begin{frame}
		\frametitle{Tecniche di I/O}
		
		\begin{enumerate}
			\item I/O programmato
			\item I/O tramite interrupt
			\item I/O con DMA
			\item I/O con canali
		\end{enumerate}
		
		\begin{center}
			\includegraphics[width=0.7\textwidth]{figures/io_techniques}
		\end{center}
		
	\end{frame}
	
	%------------------------------------------------
	\section{Obiettivi di un sistema operativo}
	
	\begin{frame}
		\frametitle{Obiettivi di un sistema operativo}
		
		\textbf{Astrazione}
		
		\begin{itemize}
			\item Alzare il livello di astrazione dei componenti hardware
			\item Semplificare l’uso del sistema
			\item Rendere più semplice programmare applicazioni
		\end{itemize}
		
		\begin{center}
			\includegraphics[width=0.7\textwidth]{figures/os_abstraction}
		\end{center}
		
	\end{frame}
	
	%------------------------------------------------
	
	\begin{frame}
		\frametitle{Virtualizzazione}
		
		\begin{itemize}
			\item Creazione di un’immagine del sistema dedicata a ogni programma
			\item Indipendenza dalla presenza di altri programmi
			\item Gestione ottimizzata delle risorse
			\item Allocazione e condivisione ordinata delle risorse
		\end{itemize}
		
	\end{frame}
	
	%------------------------------------------------
	\section{Funzioni del sistema operativo}
	
	\begin{frame}
		\frametitle{Gestione dei processi}
		
		\begin{itemize}
			\item Creazione e terminazione dei processi
			\item Sospensione e riattivazione
			\item Scheduling della CPU
			\item Sincronizzazione tra processi
			\item Gestione dei deadlock
			\item Comunicazione tra processi
		\end{itemize}
		
	\end{frame}
	
	%------------------------------------------------
	
	\begin{frame}
		\frametitle{Gestione della memoria}
		
		\begin{itemize}
			\item Allocazione e deallocazione della memoria
			\item Multiprogrammazione
			\item Caricamento dei processi
			\item Protezione della memoria
		\end{itemize}
		
	\end{frame}
	
	%------------------------------------------------
	
	\begin{frame}
		\frametitle{Gestione delle periferiche}
		
		\begin{itemize}
			\item Interazione uniforme con hardware eterogeneo
			\item Inizializzazione dei dispositivi
			\item Gestione ottimizzata dell'I/O
			\item Buffering e caching
		\end{itemize}
		
	\end{frame}
	
	%------------------------------------------------
	
	\begin{frame}
		\frametitle{Gestione del file system}
		
		\begin{itemize}
			\item File e directory
			\item Creazione e cancellazione file
			\item Lettura e scrittura
			\item Copia e ricerca
			\item Protezione e sicurezza
		\end{itemize}
		
	\end{frame}
	
	%------------------------------------------------
	\section{Avvio del sistema operativo}
	
	\begin{frame}
		\frametitle{Boot del sistema operativo}
		
		Passi principali:
		
		\begin{enumerate}
			\item Accensione
			\item Esecuzione BIOS
			\item Bootstrap loader
			\item Caricamento kernel
			\item Avvio del sistema operativo
		\end{enumerate}
		
	\end{frame}
	
	%------------------------------------------------
	
	\begin{frame}
		\frametitle{Boot loader moderni}
		
		GNU GRUB è uno dei boot loader più diffusi nei sistemi Linux.
		
		\begin{itemize}
			\item Cerca i sistemi operativi disponibili
			\item Mostra una schermata di selezione
			\item Avvia il sistema scelto
		\end{itemize}
		
		\begin{center}
			\includegraphics[width=0.6\textwidth]{figures/grub}
		\end{center}
		
	\end{frame}
	
	%------------------------------------------------
	\section{Dual Mode}
	
	\begin{frame}
		\frametitle{Dual Mode}
		
		La CPU può operare in due modalità:
		
		\begin{itemize}
			\item User mode
			\item Kernel mode
		\end{itemize}
		
		Il bit di modalità nel registro PSW indica lo stato.
		
	\end{frame}
	
	%------------------------------------------------
	\section{Interrupt}
	
	\begin{frame}
		\frametitle{Interruzioni}
		
		\begin{itemize}
			\item Interrupt hardware
			\item System call
			\item Trap ed eccezioni
		\end{itemize}
		
		\begin{center}
			\includegraphics[width=0.7\textwidth]{figures/interrupt_flow}
		\end{center}
		
	\end{frame}
	
	%------------------------------------------------
	\section{System Call}
	
	\begin{frame}
		\frametitle{Chiamate di sistema}
		
		Le system call permettono ai programmi utente
		di richiedere servizi al sistema operativo.
		
		\begin{itemize}
			\item gestione file
			\item gestione processi
			\item comunicazione
		\end{itemize}
		
	\end{frame}
	
	%------------------------------------------------
	
	\begin{frame}
		\frametitle{Esempio}
		
		Un programma C può usare:
		
		\begin{verbatim}
			printf("Hello");
		\end{verbatim}
		
		che internamente invoca la system call
		
		\begin{verbatim}
			write()
		\end{verbatim}
		
	\end{frame}
	
	%------------------------------------------------
	\section{Struttura dei sistemi operativi}
	
	\begin{frame}
		\frametitle{Strutture dei sistemi operativi}
		
		Principali architetture:
		
		\begin{itemize}
			\item sistema monolitico
			\item sistema stratificato
			\item microkernel
			\item moduli
			\item macchine virtuali
			\item client-server
		\end{itemize}
		
	\end{frame}
	
	%------------------------------------------------
	
	\begin{frame}
		\frametitle{Microkernel}
		
		Il microkernel mantiene nel kernel solo funzioni essenziali:
		
		\begin{itemize}
			\item gestione processi
			\item memoria
			\item comunicazione
		\end{itemize}
		
		Altre funzioni sono implementate come processi.
		
	\end{frame}
	
	%------------------------------------------------
	
	\begin{frame}
		\frametitle{Macchine virtuali}
		
		Una macchina virtuale crea un ambiente hardware simulato.
		
		Vantaggi:
		
		\begin{itemize}
			\item isolamento
			\item sicurezza
			\item esecuzione di più sistemi operativi
		\end{itemize}
		
	\end{frame}
	
	%------------------------------------------------
	
	\begin{frame}
		\frametitle{Conclusioni}
		
		Caratteristiche dei sistemi operativi moderni:
		
		\begin{itemize}
			\item microkernel
			\item multi-threading
			\item multi-processing
			\item sistemi distribuiti
		\end{itemize}
		
	\end{frame}
	
\end{document}